\section{Independent raw-kernel verifications}\label{sec:realizations}
\subsection{Measurable densities and hidden-state experiments}
\begin{lemma}[A measurable kernel density]\label{lem:rn}
Let $U,Y$ be standard Borel spaces, $\sigma$ a probability measure on $Y$, and $u\mapsto K_u$ a measurable finite kernel with $K_u\ll\sigma$ for every $u$. There is a jointly measurable density $f(u,y)$ representing every $K_u$ up to its own $\sigma$-null set. Consequently, sequential mixtures of such kernels have jointly measurable finite-word densities, even when hidden transitions depend on earlier reports.
\end{lemma}
\begin{proof}
Choose increasing finite partitions $\mathcal P_n$ generating the Borel sigma-field of $Y$. On a positive-mass cell $C$ set
$f_n(u,y)=K_u(C)/\sigma(C)$ and set it to zero on null cells. Each $f_n$ is jointly measurable. For fixed $u$, if $f_u$ is any Radon--Nikodym density, $f_n(u,\cdot)$ is its conditional expectation on $\mathcal P_n$. Martingale convergence gives $f_n\to f_u$ almost surely and in $L^1(\sigma)$. The liminf, set to zero where infinite, is the required joint version. The exceptional set may depend on $u$; joint measurability, not a common exceptional set for all $u$, is the assertion. Apply standard Borel disintegration to the hidden-history kernel at each step and integrate this joint density against the conditional hidden law. Iterating conditional kernels gives the product of the successive report densities. One does not condition on a hidden path as though it were independent of the reports. These version and disintegration facts agree with the standard treatment in \cite{Kallenberg,BertsekasShreve}.
\end{proof}

\begin{theorem}[Hidden-state realization]\label{thm:hidden}
Take a standard Borel hidden state, a prepared law, and marked causal transition kernels producing report, audit and boundary flag. Suppose that at each fixed action the conditional report kernel, given the hidden state and visible history, is dominated by a common $\sigma_a$. Suppose the charged joint preparation/observer-reset protocol and the all-controller return envelope \eqref{eq:UI} hold. Then this experiment verifies Theorem~\ref{thm:renewal}, without a finite-dimensional belief assumption or a filter-contraction assumption.

A complete effective verification is available for the following subclass. Apart from an independent renewal clock, the hidden space and audit alphabet are finite. The continuous part of a report is a finite union of compact cubes, with a common Lebesgue reference and finitely many visible atoms. Joint marked transition densities have uniform computable bounds on their first report derivatives and a uniform computable parameter modulus on a computably compact parameter set. The clock law and its tail envelope have certified truncations, and the preparation, density entries and compact readout loss have certified integrations. Then Theorem~\ref{thm:compile} applies, with report reconstruction error $\kappa_h\le C h$ on a mesh of diameter $h$ and a primitive model modulus obtained by integrating the stated entrywise parameter bounds.
\end{theorem}
\begin{proof}
Lemma~\ref{lem:rn} applied to the sequential conditional mixtures gives \eqref{eq:product} for every finite intervention word. The performed payoff and survival submeasures are dominated by that report law. The renewal protocol supplies iid cycles conditional on the fixed parameter, and the assumed envelope supplies precisely \eqref{eq:UI}. This proves the exact verification without invoking a realization-specific risk theorem.

For the effective subclass, keep each atom and each visible flag in its own cell. On a cube cell, replace each matrix entry density $k(y)$ by its average. A uniform derivative bound $B$ gives $|k(y)-k_C|\le B h$. Sum over the finitely many next-state, audit and flag indices and integrate over the reference volume. Half of the resulting $L^1$ bound is the required operational $\kappa_h\le C h$. The identical entrywise integration gives the parameter error $\eta(s)$, and Lemma~\ref{lem:primitive} gives
\begin{equation}\label{eq:hidden-effective}
 E_h\le\chi(\overline\mu C h),\qquad
 \Omega(s)\le\chi(\overline\mu\eta(s))+\omega_{\rm loss}(s).
\end{equation}
The clock is retained exactly by the report reconstruction. It belongs to the physical experiment, not to an uncharged observer counter. For a finite integration horizon its possible termination times below that horizon form a finite list, with the surviving tail retained as a single marked case. Finite products of computable density entries and finite probability tables have certified integrals on the cells. The uniform clock tail then bounds the omitted cycle reward and length, exactly as in the proof of Theorem~\ref{thm:compile}. This derives all of that theorem's effective inputs.

For Gaussian reports, the same construction works when the parameter family gives uniform derivative bounds on each truncation cube and a computable uniform tail bound. Keep the entire tail as a genuine cell, rather than condition it away. A coarse tail reconstruction has operational error at most its tail mass, while compact-cell errors are bounded as above. Certified Gaussian integrals and a decreasing tail bound produce $\kappa_h\to0$. A parameter change in a hidden deterministic transition is not automatically small in operational norm; that comparison must be verified independently.
\end{proof}

\begin{theorem}[A noncontractive hidden predictive geometry]\label{thm:hidden-geometry}
Let the hidden state have $d+1$ values, $d\ge1$, and a known preparation $\pi_i>0$. The first data report has reference uniform measure on $[-1,1]^d$ and likelihoods
\begin{equation}\label{eq:hidden-load}
 g_i(y)=1+\eta y_i\ (1\le i\le d),\qquad
 g_0(y)=1-\eta\sum_{i=1}^d y_i,\qquad 0<\eta<1/(2d).
\end{equation}
After this report, each next hidden transition is a known stochastic matrix $Q_a$, selected by a publicly reported innovation independent of the hidden state, under a fixed exploration. The audit target is the current hidden-state indicator vector. The cycle contains $L\ge1$ data calls, with $L$ independent of the hidden experiment and $\E(L+1)^3<\infty$. Then the actual checkpoint and online squared-prediction excess risks have matching order
\begin{equation}\label{eq:hidden-rate}
 M^{-2/d}
\end{equation}
at average risk, uniformly at every external $N\ge2$, and uniformly for $\beta\ge1/2$. All observer states, including the reset state, are counted. The matrices may include permutations and need not contract.
\end{theorem}
\begin{proof}
All likelihoods are strictly positive, integrate to one, and have constant sum $d+1$. The first posterior is
$p_i=\pi_i g_i/(\sum_j\pi_jg_j)$. It has the explicit inverse
\[
 g_i=\frac{(d+1)p_i/\pi_i}{\sum_jp_j/\pi_j},\qquad
 y_i=(g_i-1)/\eta\quad(1\le i\le d).
\]
On the compact likelihood cube this is a smooth bijection with nonsingular derivative between its $d$-dimensional interiors; the derivative and inverse derivative are uniformly bounded. The report density $\sum_i\pi_i g_i$ is bounded above and below. The first posterior law therefore has a bounded density on the simplex's affine span, and positive mass. Its conditional future-test geometry has dimension $d$: the indicator audit probabilities separate posterior states. No covariance or Fisher geometry has been postulated.

Subsequent updates are $p\mapsto pQ_a$, which are nonexpansive in $\ell^1$. The entire simplex has global covers of radius $C M^{-1/d}$ including the initial prior, and the readout into Euclidean indicator means is Lipschitz. The first likelihood update is globally Lipschitz because its denominator is uniformly positive; its initial error is zero. Corollary~\ref{cor:noncontractive} therefore supplies a uniform occupation stability bound. The first report law has total-variation continuity in its prior by integrating $\sum_i(p_i-p_i')g_i$; later public reports are independent of that prior. This independently verifies report continuity.

There is one first data report per cycle. At average risk its acquired submeasure is the first posterior law divided by $\mu=\E(L+1)$. For finite $N$, renewal epochs $s\le N-2$ each supply this same law. Their expected number exceeds $(N-2)/\mu$ by the stopped iid sum identity; it is also at least one. Thus its weight divided by $N$ is at least $1/(2\mu)$ for every $N\ge2$. At discount $\beta$ its weight is
$\beta(1-\beta)/(1-\E\beta^\tau)\ge\beta/\mu$.
For $\beta\ge1/2$ the same lower bound follows. Apply the actual small-ball inequality and the uniform global-cover recursion of Theorem~\ref{thm:geometry}. Small $M$ are absorbed into fixed positive constants. Fixed computable coefficients and a clock with certified tails also verify the effective subclass above, yielding finite executable tables with the same risk order and unchanged $M$.
\end{proof}

\subsection{Noncommuting quantum instruments}
A finite-dimensional quantum system below may be driven by a separate classical renewal clock. The clock's storage and physical cost belong to the specified experiment. They are not counted as observer memory, nor treated as costless to a simulator that must reproduce them.

\begin{theorem}[Quantum realization and operator certificates]\label{thm:quantum}
Let a finite-dimensional quantum system be prepared by a density matrix and evolve under marked completely positive instrument densities
$\mathcal I_{\theta,a}(y,z,b)$, whose sum and integral are trace preserving. The visible report contains $y,b$; the actual audit mark $z$ is not feedback. All densities are measurable with respect to common $\sigma_a$, and the joint reset protocol has a controller-uniform return envelope. Then Theorem~\ref{thm:renewal} applies to one classical $M$-state observer, including projective zeros, pure preparations and noncommuting actions.

For a complete effective realization, it suffices that the Kraus matrices and preparation matrices have computable uniform approximation bounds on a finite union of report cubes and atoms, with a certified renewal-clock tail and certified finite integrals. If $A_j(y)$ and $B_j(y)$ are corresponding Kraus approximations, the operational one-call error is bounded by one half of
\begin{equation}\label{eq:kraus}
 \sum_j\int (\norm{A_j(y)}+\norm{B_j(y)})
                       \norm{A_j(y)-B_j(y)}\,d\sigma(y).
\end{equation}
The same formula applies to parameter approximations. Report-cell averaging of the full instrument gives $\kappa_h=O(h)$ under uniform report Lipschitz bounds, and these estimates give the $E_h,\Omega$ and certified finite integrals in Theorem~\ref{thm:compile} through Lemma~\ref{lem:primitive}.
\end{theorem}
\begin{proof}
For each finite intervention word, compose the unnormalized completely positive densities and take their trace, summing unobserved audit marks. This gives a positive $L^1$ report density of total mass at most one for a stopped word. Performing the actual audit and weighting its outcomes by $\ell(d,z)$ gives a dominated payoff density. No normalization by a possibly zero report probability occurs. Preparation and reset give independent cycles conditional on the same fixed model, so the central hypotheses are verified.

For a positive input $\rho$ of trace one,
\[
 \norm{A\rho A^*-B\rho B^*}_1
 \le(\norm A+\norm B)\norm{A-B},
\]
by adding and subtracting $B\rho A^*$ and using the ideal property of the trace norm. Summation and integration give \eqref{eq:kraus}, with the classical mark retained. Cell-averaging the complete instrument, rather than averaging and then squaring its Kraus matrices, preserves complete positivity and normalization. An instrument Lipschitz bound integrates to an $O(h)$ operational error. Each finite-table response is a finite sum of traces of products of computable matrices; certified quadrature and the renewal tail give the finite risk intervals. This is an explicit operator-to-certificate implication. It does not assert that every measurable Kraus density is computable. The argument is finite-dimensional; no unbounded-operator closure is used.
\end{proof}

\begin{theorem}[Noncommuting, changing-rank acquired geometry]\label{thm:quantum-geometry}
On a $q$-level system, $q\ge2$, let $d=q^2-1$ and let $H_1,\ldots,H_d$ be a basis of traceless Hermitian matrices. Choose a symmetric cube $C\subset\R^d$ small enough that
\[
 E(y)=I+\sum_{j=1}^dy_jH_j
\]
is uniformly positive. Prepare $I/q$. With probability $1-\omega>0$, the first data instrument has report $y\in C$, uniform reference, and Kraus density $\sqrt{E(y)}$. With probability $\omega\in[0,1)$, it instead performs a tagged rank-one projective measurement $\{P_j\}_{j=1}^q$. These two tagged instruments are mixed with their stated weights.

At each data call a fixed informationally complete POVM $\{F_z\}$ is actually performed for the audit and its outcome is not feedback. Its average back-action is $\mathcal A(X)=\sum_z\sqrt{F_z}X\sqrt{F_z}$. Subsequent visible innovations select known unitaries independently of the system, so the pre-audit predictive state updates by $X\mapsto U\mathcal A(X)U^*$. Noncommuting unitaries are allowed. Let the independent data-call lifetime $L\ge1$ satisfy $\E(L+1)^3<\infty$.

For predicting the audit indicator vector under squared loss, checkpoint and online excess risks have matching order $M^{-2/(q^2-1)}$ at average risk, uniformly in $N\ge2$, and uniformly in $\beta\ge1/2$. The proof uses actual continuous mass $(1-\omega)$ and does not replace the rank-one component by a smooth density. The conclusion is for this fixed exploration, not for unrestricted quantum control.
\end{theorem}
\begin{proof}
The continuous first instrument is normalized because $\int E(y)d\sigma(y)=I$. Under $I/q$ the report law is exactly uniform, since $\tr E(y)/q=1$, and its pre-audit posterior is $X=E(y)/q$. Thus the acquired first-state law is affine full-dimensional in the $d$-dimensional trace-one space, with bounded density on its compact image. The projective component yields the pure states $P_j$ with their actual probabilities. The continuous first-state submeasure has mass $1-\omega$, not one.

For an explicit informationally complete audit, take $F_{j,\pm}=(I\pm cH_j)/(2d)$, with $c>0$ small enough for positivity. They sum to $I$. The map $f(X)=(\tr F_zX)_z$ is injective on trace-one matrices, with upper and lower Lipschitz constants on their difference space. Consequently these actually executable audit probabilities determine the predictive state, and the first continuous response law has dimension $d$. The following back-action is included in the raw marked instrument; incompatible audits are not sampled jointly.

Trace-preserving positive maps are contractive in trace norm on Hermitian differences, as follows by applying positivity and trace preservation to the positive and negative parts. Hence $X\mapsto U\mathcal A(X)U^*$ is nonexpansive. The whole density-matrix space has global trace-norm covers of radius $C M^{-1/d}$ containing $I/q$, and $f$ is Lipschitz. The first report map is regular on its continuous positive branch and defined at the projective reports; only the exact prior is used before this first update. Corollary~\ref{cor:noncontractive} gives the required stability. First report probabilities are linear in the input matrix, giving total-variation continuity; later public reports are independent of it.

The continuous first report occurs with occupation weight $(1-\omega)/\mu$ at average risk, and at least $(1-\omega)/(2\mu)$ in the finite and discounted ranges proved in Theorem~\ref{thm:hidden-geometry}. Push this actual submeasure through $f$ and apply Theorem~\ref{thm:geometry}. Its mass enters the lower constant as in \eqref{eq:masslower}. The atom and subsequent mixed-state portions remain in the upper cover and in the actual risk. The ambient trace-one space was derived here from the raw instrument and separating audit, not supplied as an unexplained predictive geometry. At $\omega=1$ the full-dimensional lower does not follow and is not asserted.
\end{proof}

\begin{corollary}[Singular references and return laws]\label{cor:singular}
The exact hidden and quantum verifications remain valid for a common Cantor report reference, a countable report law, or finitely many observed rank strata. The effective verifications remain valid when executable partitions have certified omitted mass and instrument oscillation on each retained cell. A common renewal clock with $\PP(L\ge n)\le C n^{-r}$, $r>1+p$, supplies a finite $1+p$ moment; $r>3$ supplies the noncontractive geometric criterion. All tail, continuous and stratum probabilities retain their original mass.
\end{corollary}
\begin{proof}
The measurable-density and tensor arguments use product probability measures and $L^1$ approximation, not Lebesgue density. For a discarded tail cell of mass at most $u$ uniformly in physical inputs, replacement costs at most $u$ in operational total variation; keep that cell as an actual report rather than condition on its complement. On the retained cells use their certified oscillation. For the lifetime moment, the identity expressing $\E L^s$ as a sum of successive power increments times $\PP(L\ge n)$ converges for $r>s$. Observed rank tags are part of the common reference, and the completely positive instrument estimate does not require invertible effects. Mutually singular translated Cantor families need not have one common reference, and are not covered merely by this corollary.
\end{proof}
